\documentclass[11pt]{article}
\usepackage{reusv}
\usepackage{listings}
\usepackage{enumitem}

\lstset{
  basicstyle=\ttfamily\small,
  frame=single,
  breaklines=true,
  language=Python,
  showstringspaces=false
}
\setborderthickness{0.5pt}
\setbordermargin{1cm}
\setbordercolor{black}

\slimtitle{D3. Latin Hypercube Sampling 설계}

\begin{document}
\drawborder
\thispagestyle{firstpage}

\lightertitle{D3. Latin Hypercube Sampling 설계}
\def\lighterdate{February 12, 2026}
\lighterauthor{Suwon Lee}
\lighteraddress{}{Future Mobility Control Lab, Department of Future Mobility, Kookmin University, Seoul, Republic of Korea}
\lighteremail{suwon@kookmin.ac.kr}

\begin{abstract}
본 문서는 LEO-to-LEO 연속 추력 궤도전이의 최적 추력 프로파일을 분류하기 위해,
3차원 파라미터 공간 $\mathcal{P} \subset \mathbb{R}^3$에서 초기 학습 샘플을 생성하는
Latin Hypercube Sampling (LHS) 설계를 기술한다.
LHS는 적은 샘플 수로도 각 차원의 주변분포를 균일하게 커버하여,
후속 GP 분류기의 학습 효율을 높인다.
수학적 배경, 구현 매핑, 수치 검증을 포함한다.
\end{abstract}

\noindent 대응 코드: \texttt{src/orbit\_transfer/sampling/lhs.py}


\section{목적}

본 프로젝트는 LEO-to-LEO 연속 추력 궤도전이의 최적 추력 프로파일을 분류하기 위해,
3차원 파라미터 공간 $\mathcal{P} \subset \mathbb{R}^3$에서 초기 학습 샘플을 생성해야 한다.
파라미터 공간은 다음 세 변수로 구성된다.

\begin{table}[h]
\centering
\begin{tabular}{ccccc}
\toprule
차원 $d$ & 파라미터 키 (sorted) & 기호 & 범위 & 단위 \\
\midrule
0 & \texttt{T\_normed} & $\tilde{T}$ & $[0.5,\; 5.0]$ & $T/T_0$ \\
1 & \texttt{delta\_a} & $\Delta a$ & $[-500,\; 2000]$ & km \\
2 & \texttt{delta\_i} & $\Delta i$ & $[0,\; 15]$ & deg \\
\bottomrule
\end{tabular}
\end{table}

키 정렬은 Python \texttt{sorted()} 기준으로 \texttt{['T\_normed', 'delta\_a', 'delta\_i']} 순서이며,
코드 전체에서 이 순서를 일관되게 유지한다.

이 공간에서 초기 $n = 100$개 샘플을 효율적으로 배치하기 위해 Latin Hypercube Sampling (LHS)을 사용한다.
LHS는 적은 샘플 수로도 각 차원의 주변분포(marginal distribution)를 균일하게 커버하여,
후속 GP(Gaussian Process) 분류기의 학습 효율을 높인다.


\section{수학적 배경}

\subsection{층화 추출 원리}

LHS는 층화 추출(stratified sampling)의 다차원 확장이다.

\textbf{1차원 경우.} 구간 $[0, 1]$을 $n$개의 동일 폭 부분구간으로 분할한다:
\[
I_i = \left[\frac{i}{n},\; \frac{i+1}{n}\right), \quad i = 0, 1, \dots, n-1
\]

각 부분구간에서 정확히 하나의 샘플을 추출한다:
\[
x_i = \frac{i + U_i}{n}, \quad U_i \sim \mathrm{Uniform}(0, 1)
\]

이 구조는 $n$개 샘플만으로 전 구간을 빠짐없이 커버함을 보장한다.

\textbf{다차원 확장.} $d$차원 공간 $[0,1]^d$에서 LHS를 구성하려면,
각 차원마다 독립적인 랜덤 순열 $\pi_j$를 적용한다.
$n$개 샘플로 구성된 LHS 행렬 $\mathbf{X} \in \mathbb{R}^{n \times d}$의 원소는 다음과 같다:
\[
X_{\pi_j(i),\, j} = \frac{i + U_{i,j}}{n}, \quad i = 0, \dots, n-1, \quad j = 0, \dots, d-1
\]

여기서 $\pi_j : \{0, \dots, n-1\} \to \{0, \dots, n-1\}$는 차원 $j$에 대한 랜덤 순열이고,
$U_{i,j} \sim \mathrm{Uniform}(0, 1)$는 구간 내 랜덤 섭동이다.
이 구조는 Latin square의 다차원 일반화로서,
각 행과 열에 정확히 하나의 원소가 배치되는 성질을 $d$차원으로 확장한 것이다.

\subsection{정규화 좌표와 물리 좌표 변환}

LHS는 정규화 공간 $[0,1]^d$에서 생성된다.
물리 파라미터 공간 $\mathcal{P}$와의 변환은 각 차원별 선형 매핑으로 수행한다.

\textbf{정규화 (물리 $\to$ $[0,1]$):}
\[
x_j^{\mathrm{norm}} = \frac{x_j - l_j}{h_j - l_j}
\]

\textbf{역정규화 ($[0,1]$ $\to$ 물리):}
\[
x_j = l_j + (h_j - l_j) \cdot x_j^{\mathrm{norm}}
\]

여기서 $l_j$, $h_j$는 차원 $j$의 하한과 상한이다.
두 변환은 역함수 관계를 만족한다:
\[
\mathrm{denormalize} \circ \mathrm{normalize} = \mathrm{id}
\]

\subsection{Space-Filling 성질}

LHS가 갖는 핵심 성질은 \textbf{주변 균일성(marginal uniformity)}이다.
임의의 차원 $j$에 대해, $n$개 샘플의 $j$번째 좌표를 1차원으로 사영하면,
$n$개 부분구간 각각에 정확히 하나의 점이 존재한다. 즉:
\[
|\{i : X_{i,j} \in I_k\}| = 1, \quad \forall k = 0, \dots, n-1
\]

이 성질은 어떤 차원이든 ``빈 구간'' 없이 전 범위를 커버함을 의미한다.
이는 순수 난수 샘플링에서 발생하는 클러스터링(clustering) 문제를 구조적으로 방지한다.

\subsection{LHS vs 대안 기법 비교}

\begin{table}[h]
\centering
\begin{tabular}{ccccc}
\toprule
기법 & 샘플 수 & 주변 균일성 & 분산 & 차원 확장성 \\
\midrule
순수 난수 (MC) & $n$ & 비보장 & $O(1/n)$ & 양호 \\
정규 격자 (Grid) & $n^d$ & 보장 & 0 (결정론적) & 차원의 저주 \\
LHS & $n$ & 보장 & $\leq O(1/n)$ & 양호 \\
\bottomrule
\end{tabular}
\end{table}

\textbf{순수 난수 (Monte Carlo).}
각 좌표를 독립 균일분포에서 추출한다.
샘플 수 $n$에 대해 적분 추정의 분산이 $O(1/n)$이지만,
유한 $n$에서 점들이 한곳에 밀집(clustering)되거나 빈 영역이 생길 수 있다.

\textbf{정규 격자 (Full Factorial).}
각 차원을 $m$개 수준으로 나누면 총 $m^d$개 점이 필요하다.
$d = 3$이고 차원당 10개 수준이면 $10^3 = 1000$개이다.
해석 비용이 높은 경우 비현실적이다.

\textbf{LHS.}
$n$개 샘플만으로 모든 1차원 주변분포의 균일 커버를 보장한다.
McKay et al.~\cite{McKay1979}은 LHS의 적분 추정 분산이 순수 난수 대비 항상 같거나 작음을 증명하였다.
특히 출력이 주로 주효과(main effect)에 의존하는 함수에서 분산 감소 효과가 크다.

\subsection{초기 샘플 수 결정}

본 프로젝트에서 초기 LHS 샘플 수를 $n_{\mathrm{init}} = 100$으로 설정하였다.
근거는 다음과 같다.

\begin{itemize}[nosep]
  \item \textbf{차원 대비 배수:}
    $d = 3$ 차원에서 약 $33d \approx 100$으로,
    컴퓨터 실험 설계의 경험적 지침인 $10d$를 충분히 상회한다~\cite{Loeppky2009}.
  \item \textbf{GP 학습 최소 요구량:}
    GP 분류기가 3-class 경계를 안정적으로 학습하려면,
    각 클래스에 최소 수십 개 샘플이 필요하다.
    $n = 100$은 3개 클래스에 대해 클래스당 평균 $\sim 33$개를 확보한다.
  \item \textbf{후속 적응적 샘플링과의 균형:}
    초기 샘플이 너무 적으면 GP 예측이 불안정하고,
    너무 많으면 적응적 샘플링의 이점이 줄어든다.
    최대 샘플 수 $n_{\max} = 500$ 대비 20\%를 초기 탐색에 할당하는 것은 적절한 비율이다.
\end{itemize}


\section{구현 매핑}

\subsection{함수--수학 대응표}

\begin{table}[h]
\centering
\begin{tabular}{llll}
\toprule
수학 표현 & Python 함수 & 파일 위치 & 비고 \\
\midrule
LHS 알고리즘 (섹션 2.1) & \texttt{latin\_hypercube\_sample()} & \texttt{lhs.py:8--41} & 순열 + 구간 내 균일 랜덤 \\
$\mathcal{P} \to [0,1]^d$ 정규화 & \texttt{normalize\_params()} & \texttt{lhs.py:44--53} & 차원별 선형 스케일링 \\
$[0,1]^d \to \mathcal{P}$ 역정규화 & \texttt{denormalize\_params()} & \texttt{lhs.py:56--65} & 차원별 역선형 스케일링 \\
\bottomrule
\end{tabular}
\end{table}

\subsection{알고리즘 흐름}

\texttt{latin\_hypercube\_sample(n\_samples, param\_ranges, seed)} 함수의 실행 흐름은 다음과 같다.

\begin{enumerate}[nosep]
  \item \textbf{초기화:}
    \texttt{param\_ranges}가 \texttt{None}이면 \texttt{config.PARAM\_RANGES}를 사용한다.
    키를 \texttt{sorted()}로 정렬하여 차원 순서를 결정한다.
  \item \textbf{정규화 좌표 생성 (lines 29--33):}
    \begin{itemize}[nosep]
      \item 각 차원 $j = 0, \dots, d-1$에 대해:
        \begin{itemize}[nosep]
          \item $\pi_j \leftarrow$ \texttt{rng.permutation(n\_samples)} --- 랜덤 순열 생성
          \item 각 구간 $i = 0, \dots, n-1$에 대해:
            \texttt{samples\_normed[$\pi_j(i)$, $j$] = ($i$ + $U$) / $n$} --- $U \sim \mathrm{Uniform}(0,1)$
        \end{itemize}
    \end{itemize}
  \item \textbf{물리 좌표 변환 (lines 36--39):}
    각 차원 $j$에 대해 \texttt{samples[:, $j$] = lo + (hi - lo) * samples\_normed[:, $j$]}
  \item \textbf{반환:}
    \texttt{(samples, samples\_normed)} --- 물리 좌표와 정규화 좌표를 함께 반환
\end{enumerate}

\texttt{normalize\_params}와 \texttt{denormalize\_params}는 임의 shape의 배열을 지원한다.
마지막 축(axis $= -1$)의 크기가 $d$이면 되며,
\texttt{...} 인덱싱으로 브로드캐스팅을 처리한다.

\subsection{난수 생성기}

\texttt{numpy.random.default\_rng(seed)}를 사용하여 PCG-64 기반 난수 생성기를 초기화한다.
\texttt{seed}를 명시하면 결과가 재현 가능하다.
이는 실험 재현성 확보에 필수적이다.


\section{수치 검증}

대응 테스트: \texttt{tests/test\_sampling.py::TestLHS}

\subsection{형상 검증}

\begin{lstlisting}
def test_shape(self):
    samples, normed = latin_hypercube_sample(50, seed=42)
    assert samples.shape == (50, 3)
    assert normed.shape == (50, 3)
\end{lstlisting}

$n = 50$개 샘플, $d = 3$ 차원에서 출력 행렬의 shape이 $(50, 3)$인지 확인한다.

\subsection{범위 검증}

\begin{lstlisting}
def test_range(self):
    _, normed = latin_hypercube_sample(100, seed=42)
    assert np.all(normed >= 0) and np.all(normed <= 1)
\end{lstlisting}

정규화 좌표가 $[0, 1]$ 범위 내에 있는지 확인한다.
LHS 구성상 $x_{i,j} = (i + U_{i,j})/n$이므로,
$U_{i,j} \in (0,1)$에서 $x_{i,j} \in (0/n,\; (n-1+1)/n) = (0, 1)$이 보장된다.

\subsection{주변 균일성 검증}

\begin{lstlisting}
def test_uniformity(self):
    _, normed = latin_hypercube_sample(1000, seed=42)
    for d in range(3):
        hist, _ = np.histogram(normed[:, d], bins=10)
        assert np.all(hist > 70)
        assert np.all(hist < 130)
\end{lstlisting}

$n = 1000$개 샘플의 각 차원을 10개 bin으로 히스토그램화한다.
이상적 LHS에서 각 bin에 정확히 $1000/10 = 100$개가 배치되어야 한다.
허용 범위 $[70, 130]$ ($\pm 30\%$)으로 검증한다.
LHS의 층화 구조 덕분에 실제로는 이보다 훨씬 균일한 분포가 나타난다.

\subsection{Round-Trip 일관성}

\begin{lstlisting}
def test_normalize_denormalize_roundtrip(self):
    samples, normed = latin_hypercube_sample(20, seed=0)
    recovered = denormalize_params(normalize_params(samples))
    np.testing.assert_allclose(recovered, samples, atol=1e-12)
\end{lstlisting}

물리 좌표 $\to$ 정규화 $\to$ 역정규화 경로에서 원본이 복원되는지 확인한다.
허용 오차 $10^{-12}$는 부동소수점 산술의 기계 정밀도 수준이다.

\subsection{키 정렬 순서}

\texttt{sorted(PARAM\_RANGES.keys())}의 결과는 \texttt{['T\_normed', 'delta\_a', 'delta\_i']}이다.
\texttt{latin\_hypercube\_sample}, \texttt{normalize\_params}, \texttt{denormalize\_params} 세 함수 모두
동일한 \texttt{sorted()} 호출을 사용하므로,
차원 인덱스 $d = 0, 1, 2$가 항상 같은 물리 파라미터에 대응된다.


\bibliographystyle{IEEEtran}
\bibliography{references}

\end{document}
